% ===============================================================
%  Chapter 5 -- The MSI Family: MSI, MESI and MOESI      (Part I)
% ===============================================================
\chapter{The MSI Family: MSI, MESI and MOESI}
\label{ch:msi-family}
\index{MESI}\index{MOESI}\index{MSI}%

\begin{figure}[htbp]
\centering
\adjustbox{max width=\linewidth}{%
\begin{forest}
  chapter map,
  [Chapter 5
    [Events {\&} messages [Processor events] [Bus messages]]
    [MSI [Three states] [Two state diagrams]]
    [MESI [Adds Exclusive] [Silent upgrade]]
    [MOESI [Adds Owned] [Read-miss decision tree]]
    [A worked trace [Same accesses in all three] [The family at a glance]]
  ]
\end{forest}}
\caption[]{Structure of Chapter~5. Each protocol is given
in the same order (states, transitions on processor events, transitions on
messages from other caches) so the three can be compared line by line.}
\label{fig:ch5-map}
\end{figure}

This chapter describes what each of the three protocols does. Which one costs
less is the subject of Parts~II and~III.

% ---------------------------------------------------------------
\section{Events and Messages}
\label{sec:fam-events}
\sectionmap{[Vocabulary [Processor events: read write evict]
  [Bus messages: readmiss writemiss write] [Responses: data and write back]
  [Hits change nothing]]}
\index{write-invalidate}\index{Sarangi textbook}%

The three protocols react to the same events and use the message names of
Sarangi's write-invalidate protocol~\cite[\S12.4.6]{sarangi2025basic}
(\cref{tab:events}). Read hits in any valid state and write hits in \stM\ change
no state and send no message, so the diagrams leave them out.

\begin{table}[htbp]
\centering
\caption{Events and messages used by the MSI family.}
\label{tab:events}
\small
\begin{tabular}{@{}l l p{7.6cm}@{}}
\toprule
Kind & Name & Meaning \\
\midrule
Processor event & read / write & the local core loads or stores the block \\
 & evict & the local cache removes the block to make room \\
\midrule
Bus message & \textit{readmiss} & another cache wants to read a block it does not hold \\
 & \textit{writemiss} & another cache wants to write a block it does not hold \\
 & \textit{write} & another cache holds the block and is about to write it
   (an upgrade); other copies must be invalidated \\
\midrule
Response & \textit{send data} & a cache supplies the block to the requester \\
 & \textit{write back} & a cache writes the block to the lower level \\
\bottomrule
\end{tabular}
\end{table}

% ---------------------------------------------------------------
\section{MSI: Three States}
\label{sec:fam-msi}
\sectionmap{[MSI [M: only copy and modified] [S: read-only and clean]
  [I: not present] [One writer or many readers]]}
\index{MSI}\index{write-invalidate}\index{Sarangi textbook}%

The basic write-invalidate protocol, \gls{msi}\index{MSI}, has three states per
block~\cite[\S12.4.6]{sarangi2025basic}:

\begin{description}
  \item[\stM\ (Modified).] This cache may write the block. No other cache
    holds a valid copy. The data differs from memory.
  \item[\stS\ (Shared).] This cache may read but not write the block. Other
    caches may hold \stS\ copies too. The data matches memory.
  \item[\stI\ (Invalid).] The cache holds no valid copy.
\end{description}

\Cref{fig:msi} redraws Sarangi's two diagrams: transitions caused by the local
processor (his Fig.~12.13) and transitions caused by messages from other caches
(his Fig.~12.14). The red transition in part~(b) comes up again and again in this
report. When another cache reads a block held in \stM, this cache sends the data,
writes the block back to the lower level, and moves to \stS.

\tikzset{
  fsm/.style={font=\sffamily\small, >=Stealth, node distance=2cm},
  st/.style={draw, thick, circle, minimum size=1.05cm, font=\sffamily\bfseries},
  stB/.style={st, fill=mesiBlue!14, draw=mesiBlue},
  stO/.style={st, fill=moesiOrange!14, draw=moesiOrange},
  stT/.style={st, fill=tealC!16, draw=tealC},
  tr/.style={->, thick, slateC},
  bus/.style={->, thick, dashed, plumC},
  key/.style={->, very thick, alertRed},
  good/.style={->, very thick, memGreen},
  lb/.style={font=\sffamily\tiny, fill=white, inner sep=1pt, align=center},
}

\begin{figure}[htbp]
\centering
\begin{bigfig}
\begin{tikzpicture}[fsm]
  \node[font=\sffamily\bfseries\small] at (2.4,4.2) {(a) Processor events};
  \node[stT] (I) at (0,0) {\stI};
  \node[stT] (S) at (4.8,0) {\stS};
  \node[stT] (M) at (2.4,3.0) {\stM};
  \draw[tr] (I) to[bend right=18] node[lb, below]{read / readmiss} (S);
  \draw[tr] (S) to[bend right=18] node[lb, above]{evict / ---} (I);
  \draw[tr] (I) to[bend left=12] node[lb, left=2pt]{write /\\writemiss} (M);
  \draw[tr] (M) to[bend left=12] node[lb, right=1pt]{evict /\\write back} (I);
  \draw[tr] (S) -- node[lb, right=3pt]{write /\\write} (M);
  \begin{scope}[xshift=7.4cm]
  \node[font=\sffamily\bfseries\small] at (2.4,4.2) {(b) Messages from other caches};
  \node[stT] (I2) at (0,0) {\stI};
  \node[stT] (S2) at (4.8,0) {\stS};
  \node[stT] (M2) at (2.4,3.0) {\stM};
  \draw[bus] (S2) -- node[lb, below]{write, writemiss / ---} (I2);
  \draw[bus] (M2) -- node[lb, left=2pt]{writemiss /\\send data} (I2);
  \draw[key] (M2) -- node[lb, right=3pt, text=alertRed]{readmiss /\\send data\\+ \textbf{write back}} (S2);
  \end{scope}
\end{tikzpicture}
\end{bigfig}
\caption[MSI state diagrams]{The MSI protocol, redrawn from Sarangi's
Figs.~12.13 and~12.14~\cite[\S12.4.6]{sarangi2025basic}. Solid arrows are caused
by the local processor and dashed purple arrows by messages from other caches. The
red transition writes the block back when another cache reads it.}
\label{fig:msi}
\end{figure}

% ---------------------------------------------------------------
\section{MESI: Adding an Exclusive State}
\label{sec:fam-mesi}
\sectionmap{[MESI [Read of unshared data enters E] [E is clean and the only copy]
  [E to M needs no message] [Origin: Illinois protocol 1984]]}
\index{MESI}\index{MSI}\index{Exclusive state}\index{Illinois protocol}\index{Sarangi textbook}%

When a core reads a block that no other cache holds, such as its own stack
data, MSI still puts the block in \stS\ because it cannot tell that this is the
only copy. When the core then writes the block, the cache has to send a
\textit{write} message to invalidate copies that do not exist.
\Gls{mesi}\index{MESI} adds a fourth state for this case~\cite{papamarcos1984illinois}:

\begin{description}
  \item[\stE\ (Exclusive).] No other cache has a copy and the data matches
    memory. The cache may read the block and may move it to \stM\ on a write
    without sending any message, since no other cache needs to know.
\end{description}

On a read miss the other caches report whether any of them holds the block, for
example by asserting a ``shared'' line on the bus. If none does, the block enters
\stE\ instead of \stS. Sarangi introduces this state and leaves ``implementing
the MESI protocol\,\ldots as an exercise for the
reader''~\cite[\S12.4.6, p.~603]{sarangi2025basic}; \cref{fig:mesi} is our
solution.

\begin{figure}[htbp]
\centering
\adjustbox{max width=\linewidth}{%
\begin{tikzpicture}[fsm]
  \node[font=\sffamily\bfseries\small] at (2.9,4.1) {(a) Processor events};
  \node[stB] (I) at (0,1.5) {\stI};
  \node[stB] (E) at (2.9,3.2) {\stE};
  \node[stB] (S) at (2.9,-0.2) {\stS};
  \node[stB] (M) at (5.8,1.5) {\stM};
  \draw[tr] (I) -- node[lb, above left=-1pt]{read, no copies /\\readmiss} (E);
  \draw[tr] (I) -- node[lb, below left=-1pt]{read, copies exist /\\readmiss} (S);
  \draw[tr] (I) to[bend left=9] node[lb, above, pos=0.62]{write / writemiss} (M);
  \draw[tr] (M) to[bend left=9] node[lb, below, pos=0.62]{evict / write back} (I);
  \draw[good] (E) -- node[lb, above right=-1pt, text=memGreen]{write / ---\\\textbf{(silent)}} (M);
  \draw[tr] (S) -- node[lb, below right=-1pt]{write / write} (M);
  \node[lb, text=coreGrey, anchor=west] at (-0.4,-1.0) {Evict from \stE\ or \stS: silent, to \stI.};
  \begin{scope}[xshift=7.7cm, yshift=0.4cm]
  \node[font=\sffamily\bfseries\small] at (2.8,3.8) {(b) Messages from other caches};
  \node[stB] (E2) at (0,2.4) {\stE};
  \node[stB] (S2) at (2.8,2.4) {\stS};
  \node[stB] (M2) at (5.6,2.4) {\stM};
  \node[stB] (I2) at (2.8,-0.4) {\stI};
  \draw[bus] (E2) -- node[lb, above=3pt]{readmiss /\\send data} (S2);
  \draw[key] (M2) -- node[lb, above=3pt, text=alertRed]{readmiss / send data\\+ \textbf{write back}} (S2);
  \draw[bus] (E2) -- node[lb, left=2pt, pos=0.45]{writemiss /\\send data} (I2);
  \draw[bus] (S2) -- node[lb, pos=0.45]{write,\\writemiss / ---} (I2);
  \draw[bus] (M2) -- node[lb, right=2pt, pos=0.45]{writemiss /\\send data} (I2);
  \end{scope}
\end{tikzpicture}}
\caption[MESI state diagrams]{The MESI protocol. The green transition is what
MESI adds: a block in \stE\ moves to \stM\ silently. The red transition is the
same as in MSI, so a remote read of an \stM\ block still causes a write back.}
\label{fig:mesi}
\end{figure}

% ---------------------------------------------------------------
\section{MOESI: Adding an Owned State}
\label{sec:fam-moesi}
\label{sec:fam-readmiss}
\sectionmap{[MOESI [O: modified but shared] [Owner supplies data]
  [Owner is responsible for write back] [Origin: Futurebus 1986]]}
\index{MESI}\index{MOESI}\index{Owned state}\index{snooping protocol}\index{Futurebus}\index{owner}%

\Gls{moesi}\index{MOESI} adds a fifth state to MESI~\cite{sweazey1986moesi}:

\begin{description}
  \item[\stO\ (Owned).] The copy is modified, so memory is stale, but other
    caches may hold the same data in \stS. The owning cache answers every request
    for the block and must write it back to memory when it evicts it.
\end{description}

The only change from MESI is the red transition of \cref{fig:mesi}. In MOESI, when
another cache reads a block held in \stM, the holder sends the data and moves to
\stO\ without writing back (\cref{fig:moesi}). An \stS\ copy may then differ from
memory, because the latest data is in the \gls{owner}\index{owner}, and the owner
must write the block back when it evicts it, as an \stM\ block must.
\Cref{fig:readmiss-tree} shows how a snooping system serves a read miss under both
protocols. They agree on every branch except the last.

\begin{figure}[htbp]
\centering
\adjustbox{max width=\linewidth}{%
\begin{forest}
  decision tree,
  [Read miss: does another\\cache hold the block?
    [Requester enters E\\data from memory, el=no]
    [Is that copy\\in M or O?, el=yes
      [Requester enters S\\any E copy drops to S, el=no]
      [Which protocol?, el=yes
        [MESI: owner sends data\\writes back and enters S, el=MESI]
        [MOESI: owner sends data\\enters or stays in O\\no write back, el=MOESI]]]]
\end{forest}}
\caption[Decision tree: serving a read miss]{How a read miss is served in MESI
and MOESI. The protocols differ only in the bottom row, where a cache holding
modified data is read by another cache.}
\label{fig:readmiss-tree}
\end{figure}


\begin{figure}[htbp]
\centering
\begin{bigfig}
\begin{tikzpicture}[fsm]
  \node[font=\sffamily\bfseries\small, anchor=west] at (-0.5,3.2) {(a) Processor events: the two transitions new in MOESI};
  \node[stO] (M) at (1.9,2.0) {\stM};
  \node[stO] (O) at (3.8,0) {\stO};
  \node[stO] (I) at (0,0) {\stI};
  \draw[tr] (O) -- node[lb, right=2pt]{write /\\write} (M);
  \draw[tr] (O) -- node[lb, below]{evict / write back} (I);
  \node[lb, text=coreGrey, align=left, anchor=west] at (5.0,1.0)
       {All other processor-side transitions are\\exactly as in MESI (\cref{fig:mesi}a).};
  \begin{scope}[yshift=-5.6cm]
  \node[font=\sffamily\bfseries\small, anchor=west] at (-0.5,3.9) {(b) Messages from other caches};
  \node[stO] (E2) at (0,2.4) {\stE};
  \node[stO] (S2) at (3.0,2.4) {\stS};
  \node[stO] (M2) at (6.0,2.4) {\stM};
  \node[stO] (O2) at (9.0,2.4) {\stO};
  \node[stO] (I2) at (4.5,-0.6) {\stI};
  \draw[bus] (E2) -- node[lb, above=3pt]{readmiss /\\send data} (S2);
  \draw[good] (M2) -- node[lb, above=3pt, text=memGreen]{readmiss / send data\\\textbf{no write back}} (O2);
  \draw[bus] (O2) to[loop right, looseness=5] node[lb, right]{readmiss /\\send data} (O2);
  \draw[bus] (E2) -- node[lb, left=3pt, pos=0.5]{writemiss / send data} (I2);
  \draw[bus] (S2) -- node[lb, pos=0.45]{write,\\writemiss / ---} (I2);
  \draw[bus] (M2) -- node[lb, pos=0.45]{writemiss /\\send data} (I2);
  \draw[bus] (O2) -- node[lb, right=3pt, pos=0.5]{write, writemiss / send data if asked} (I2);
  \end{scope}
\end{tikzpicture}
\end{bigfig}
\caption[MOESI state diagrams]{The MOESI protocol. Part~(a) shows only the two
processor-side transitions that MOESI adds. In part~(b) the green transition
replaces MESI's red one: the cache with modified data supplies it and moves to
\stO\ without writing back, and then answers every later read miss itself (the
loop on \stO).}
\label{fig:moesi}
\end{figure}

Some combinations of states in different caches would break coherence, so each
protocol allows only certain pairs for the same block. \Cref{tab:compat} gives the
rule for MOESI. The one new pair is $(\stO,\stS)$, which is modified data that is
also shared.

\begin{table}[htbp]
\centering
\caption{Which states two caches may hold for the same block at the same time
in MOESI ($\checkmark$ = allowed). Removing the \stO\ row and column gives the
rule for MESI.}
\label{tab:compat}
\small
\begin{tabular}{@{}c|ccccc@{}}
 & \stM & \stO & \stE & \stS & \stI \\
\hline
\stM & $\times$ & $\times$ & $\times$ & $\times$ & \checkmark \\
\stO & $\times$ & $\times$ & $\times$ & \checkmark & \checkmark \\
\stE & $\times$ & $\times$ & $\times$ & $\times$ & \checkmark \\
\stS & $\times$ & \checkmark & $\times$ & \checkmark & \checkmark \\
\stI & \checkmark & \checkmark & \checkmark & \checkmark & \checkmark \\
\end{tabular}
\end{table}

% ---------------------------------------------------------------
\section{A Worked Trace}
\label{sec:fam-trace}
\sectionmap{[Worked trace [Three cores and one block] [Six accesses]
  [States in MSI MESI MOESI] [E appears at step 1] [O appears at step 3]]}
\index{MESI}\index{MOESI}\index{MSI}\index{Owned state}\index{Exclusive state}\index{owner}%

\Cref{tab:trace} follows one block through six accesses by three cores under
each protocol. The block starts in no cache.

\begin{table}[htbp]
\centering
\caption{One block, six accesses, three protocols. Cells give the states in
P0 / P1 / P2 after each step.}
\label{tab:trace}
\small
\begin{tabular}{@{}c l c c c@{}}
\toprule
Step & Access & MSI & MESI & MOESI \\
\midrule
1 & P0 reads  & S / I / I & \textbf{E} / I / I & \textbf{E} / I / I \\
2 & P0 writes & M / I / I & M / I / I & M / I / I \\
3 & P1 reads  & S / S / I & S / S / I & \textbf{O} / S / I \\
4 & P2 reads  & S / S / S & S / S / S & \textbf{O} / S / S \\
5 & P1 writes & I / M / I & I / M / I & I / M / I \\
6 & P0 reads  & S / S / I & S / S / I & S / \textbf{O} / I \\
\bottomrule
\end{tabular}
\end{table}

At step~1 MESI and MOESI put the block in \stE, since no other cache holds it;
MSI has no such state. At step~2 the move to \stM\ is silent in MESI and MOESI
and costs a bus \textit{write} in MSI. At steps~3, 4 and~6 MOESI keeps the
modified data in an owner (\stO), while MSI and MESI fall back to all-\stS. At
step~5 ownership passes to the new writer, and P0's \stO\ copy is invalidated
because P1's write makes it out of date.

\label{sec:fam-glance}
\Cref{tab:family} describes every state of the family by four properties and
shows which protocol uses it.

\begin{table}[htbp]
\centering
\caption{Every state of the MSI family.}
\label{tab:family}
\small
\begin{tabular}{@{}c c c c c c c c@{}}
\toprule
State & Valid & Differs from & Only & Responsible for & \multicolumn{3}{c}{Used in} \\
\cmidrule(l){6-8}
 & & memory & copy & write back & MSI & MESI & MOESI \\
\midrule
\stM & yes & yes & yes & yes & \checkmark & \checkmark & \checkmark \\
\stO & yes & yes & no  & yes &            &            & \checkmark \\
\stE & yes & no  & yes & no  &            & \checkmark & \checkmark \\
\stS & yes & no$^{*}$ & no & no & \checkmark & \checkmark & \checkmark \\
\stI & no  & --- & --- & no  & \checkmark & \checkmark & \checkmark \\
\midrule
\multicolumn{5}{@{}l}{Bits of state per block} & 2 & 2 & 3 \\
\bottomrule
\multicolumn{8}{@{}l}{\footnotesize $^{*}$ In MOESI, an \stS\ copy may differ
from memory when another cache holds the block in \stO.}
\end{tabular}
\end{table}

MESI is therefore MSI with the clean, unshared case taken out of \stS\ and given
its own state \stE, and MOESI is MESI with a dirty, shared case added as \stO.
